\chapter{Testing and Evaluation}
\label{ch:testing}

This chapter reports how the system was tested and, more importantly, what the
trained models actually do on real data. It opens with a short account of the
software test suite, then spends most of its length on the experiments: the
datasets, the detection results, the quality-classification results across the
model variants, and the measured latency of the running prototype. It closes with
the limitations the numbers expose.

\section{Testing Strategy}

The backend carries most of the software tests because it is where a wrong change
does the most damage. The suite is layered. Unit tests check single functions in
isolation with the database, the object store, and the model backends replaced by
fakes, so they run in milliseconds and cover the pure logic: request and response
validation, the batch state transitions, and the authentication token paths.
Integration tests run against a real PostgreSQL instance started for the test run
and torn down afterwards, so repository code is checked against the database
engine it actually uses rather than a mock that might disagree with real SQL.
End-to-end tests drive the assembled HTTP application and assert on status codes
and response bodies, exercising routing, authentication, and serialization
together. On top of these sits a smoke test that boots the full stack, registers
a model, submits an analyze request, and waits for the batch to finish, which is
the closest check to a real deployment.

The web client has component tests, including one that switches the interface to
Arabic and asserts that the layout flips to right-to-left and the choice
persists. MultiSeedGen has two checks that matter for a data generator: a
golden-output test that pins the exact bytes of a generated dataset so any change
to the output is caught, and a determinism test that asserts the same
configuration and random seed always produce identical output regardless of how
many worker processes run.

Table~\ref{tab:ch6-cases} lists representative cases from across the suite. The
result column reports pass where the behaviour matches the expected contract.

\begin{singlespace}
\begin{longtable}{@{}p{0.05\textwidth}p{0.13\textwidth}p{0.40\textwidth}p{0.27\textwidth}p{0.06\textwidth}@{}}
\caption{Representative software test cases.}
\label{tab:ch6-cases}\\
\toprule
\textbf{ID} & \textbf{Area} & \textbf{Scenario} & \textbf{Expected} & \textbf{Result} \\
\midrule
\endfirsthead
\toprule
\textbf{ID} & \textbf{Area} & \textbf{Scenario} & \textbf{Expected} & \textbf{Result} \\
\midrule
\endhead
\bottomrule
\endfoot
T-01 & Auth & Login with valid credentials. & An access token and a refresh token are issued. & Pass \\
\addlinespace
T-02 & Auth & A used refresh token is presented again. & The reused token is rejected and the session is invalidated. & Pass \\
\addlinespace
T-03 & Analyze & A batch of valid images is uploaded. & The batch is created and one analysis task is queued per image. & Pass \\
\addlinespace
T-04 & Analyze & A file that is not a real image is uploaded. & Validation fails and no batch is stored. & Pass \\
\addlinespace
T-05 & Pipeline & Detection succeeds but classification fails on one image. & Detections are kept and the batch finishes partial, not failed. & Pass \\
\addlinespace
T-06 & Generator & MultiSeedGen runs twice with the same seed and different worker counts. & Both runs produce byte-identical output. & Pass \\
\addlinespace
T-07 & Generator & A train/validation split is built from the seed pool. & No source seed appears in both partitions. & Pass \\
\end{longtable}
\end{singlespace}

\section{Datasets}
\label{sec:ch6-datasets}

The models were trained and evaluated on data drawn from four sources, chosen so
the system would see more than one capture condition. Public seed datasets are
usually shot in a controlled lab, and a model trained only on that kind of image
tends to fail on a photo taken with a phone under ordinary light. Mixing sources
was the main defence against that.

The first source is the NAGAR maize dataset, built on the IIIT-Hyderabad corn
seed benchmark~\cite{naagar2025corn}. We used the balanced variant and removed
noisy labels to improve quality. For the classifier the data is reduced to two
classes: good, which is healthy pure corn, and bad, which combines the broken,
discoloured, and silkcut defects.

The second source is a maize dataset from Roboflow with 26{,}350
images~\cite{roboflow2025seedquality}. Unlike the lab-controlled NAGAR set, it
spans a wide range of lighting and texture and is already split into binary good
and bad classes. Combining it with NAGAR produced what we call the hybrid
dataset.

The third source is a smaller but densely annotated detection set of 600
high-density images used to train and stress-test the localization
models~\cite{roboflow2025maize}. Every seed carries a bounding box, and the boxes
are labelled across seven maize classes: broken (2{,}927), damage (3{,}766),
fungus (1{,}510), immature (1{,}721), shriveled (1{,}618), weeveled (1{,}835),
and healthy (1{,}706).

The fourth source is a coffee parchment dataset, the dried endocarp of the coffee
bean, with good and bad classes. It was added to check whether the design
generalizes to a second crop with a different shape and spatial arrangement.

\section{Object Detection Results}
\label{sec:ch6-detection}

The first stage of the pipeline locates each seed. We evaluated Faster
R-CNN~\cite{fasterrcnn2015} in four configurations, with and without a CBAM
attention block~\cite{cbam2018}, on maize alone and on maize and coffee together.
Detection quality is reported as mean average precision at three intersection
thresholds: mAP@50, the stricter mAP@75, and the COCO average over 0.50 to 0.95.
Table~\ref{tab:ch6-frcnn} collects the results.

\begin{singlespace}
\begin{longtable}{@{}p{0.40\textwidth}p{0.13\textwidth}p{0.13\textwidth}p{0.13\textwidth}p{0.10\textwidth}@{}}
\caption{Faster R-CNN detection results across configurations.}
\label{tab:ch6-frcnn}\\
\toprule
\textbf{Configuration} & \textbf{mAP@50} & \textbf{mAP@75} & \textbf{mAP .5:.95} & \textbf{Time (s)} \\
\midrule
\endfirsthead
\toprule
\textbf{Configuration} & \textbf{mAP@50} & \textbf{mAP@75} & \textbf{mAP .5:.95} & \textbf{Time (s)} \\
\midrule
\endhead
\bottomrule
\endfoot
Faster R-CNN, maize only          & 0.9780 & 0.6574 & 0.5903 & 0.1511 \\
\addlinespace
Faster R-CNN + CBAM, maize only   & 0.9768 & 0.6576 & 0.5867 & 0.0621 \\
\addlinespace
Faster R-CNN, maize \& coffee     & 0.9838 & 0.8015 & 0.7038 & 0.1038 \\
\addlinespace
Faster R-CNN + CBAM, maize \& coffee & 0.9835 & 0.8020 & 0.7211 & 0.1106 \\
\end{longtable}
\end{singlespace}

Two patterns stand out. On the single-class maize task, adding CBAM did not help;
the standard and attention models score within a thousandth of each other on
mAP@50, so the extra complexity bought nothing on an easy, single-class set. On
the harder two-crop task the picture changes. Moving from maize-only to mixed
maize-and-coffee data lifted every metric, and the jump is largest on the strict
thresholds: mAP@75 rises from about 0.66 to about 0.80 and the COCO average from
0.59 to 0.70. Adding more variety to the training data helped the detector more
than any architectural change did. On that mixed set CBAM gave a small but real
gain on the strictest metric, nudging the COCO average from 0.7038 to 0.7211,
which fits the idea that attention earns its cost on harder, more varied inputs.

For deployments where speed matters more than the last point of accuracy, we also
trained a one-stage YOLOv8 detector~\cite{redmon2016yolo,yolov8}. It reaches an
mAP@50 of 0.941 with a precision of 0.939 and a recall of 0.941, at roughly
5.5~ms per image, which is far faster than Faster R-CNN and is why the prototype
offers it as a fast mode.

\section{Seed Quality Classification Results}
\label{sec:ch6-classification}

The second stage takes each detected seed and labels it good or bad. The
classifier is a ResNet-18~\cite{resnet2015} that we extended in four steps to
measure the effect of each change: V1 is the plain baseline, V2 adds a CBAM
attention block~\cite{cbam2018}, V3 adds global max pooling alongside the usual
average pooling, and V4 adds a stride change in the final stage. We trained this
family twice, once on the original NAGAR data and once on the hybrid set that
adds the Roboflow images. Table~\ref{tab:ch6-clf-summary} shows the headline
comparison between the baseline on each.

\begin{singlespace}
\begin{longtable}{@{}p{0.26\textwidth}p{0.22\textwidth}p{0.11\textwidth}p{0.11\textwidth}p{0.12\textwidth}@{}}
\caption{Baseline classifier on the NAGAR and hybrid datasets.}
\label{tab:ch6-clf-summary}\\
\toprule
\textbf{Model} & \textbf{Training data} & \textbf{Acc.} & \textbf{F1} & \textbf{Time} \\
\midrule
\endfirsthead
\toprule
\textbf{Model} & \textbf{Training data} & \textbf{Acc.} & \textbf{F1} & \textbf{Time} \\
\midrule
\endhead
\bottomrule
\endfoot
Baseline ResNet-18 & NAGAR only        & 96.19\% & 0.9626 & 70.05\,s \\
\addlinespace
Hybrid ResNet-18   & NAGAR + Roboflow   & 97.39\% & 0.9745 & 143.50\,s \\
\end{longtable}
\end{singlespace}

\subsection{Validation performance}

On held-out validation data the four variants all score highly, which is expected
because validation data comes from the same distribution as training.
Tables~\ref{tab:ch6-clf-nagar-val} and \ref{tab:ch6-clf-hybrid-val} report the
variants on the NAGAR and hybrid sets.

\begin{singlespace}
\begin{longtable}{@{}p{0.30\textwidth}p{0.12\textwidth}p{0.12\textwidth}p{0.12\textwidth}p{0.12\textwidth}@{}}
\caption{Classifier variants on the NAGAR dataset (validation).}
\label{tab:ch6-clf-nagar-val}\\
\toprule
\textbf{Variant} & \textbf{F1} & \textbf{Combined} & \textbf{Precision} & \textbf{Recall} \\
\midrule
\endfirsthead
\toprule
\textbf{Variant} & \textbf{F1} & \textbf{Combined} & \textbf{Precision} & \textbf{Recall} \\
\midrule
\endhead
\bottomrule
\endfoot
V1 Baseline ResNet-18    & 0.9626 & 0.9623 & 0.9768 & 0.9489 \\
\addlinespace
V2 + CBAM                & 0.9760 & 0.9757 & 0.9869 & 0.9653 \\
\addlinespace
V3 + CBAM + GMP          & 0.9758 & 0.9755 & 0.9821 & 0.9696 \\
\addlinespace
V4 + CBAM + GMP + Stride & 0.9734 & 0.9730 & 0.9797 & 0.9671 \\
\end{longtable}
\end{singlespace}

\begin{singlespace}
\begin{longtable}{@{}p{0.30\textwidth}p{0.12\textwidth}p{0.12\textwidth}p{0.12\textwidth}p{0.12\textwidth}@{}}
\caption{Classifier variants on the hybrid dataset (validation).}
\label{tab:ch6-clf-hybrid-val}\\
\toprule
\textbf{Variant} & \textbf{F1} & \textbf{Combined} & \textbf{Precision} & \textbf{Recall} \\
\midrule
\endfirsthead
\toprule
\textbf{Variant} & \textbf{F1} & \textbf{Combined} & \textbf{Precision} & \textbf{Recall} \\
\midrule
\endhead
\bottomrule
\endfoot
V1 Baseline ResNet-18    & 0.9779 & 0.9776 & 0.9846 & 0.9714 \\
\addlinespace
V2 + CBAM                & 0.9789 & 0.9786 & 0.9846 & 0.9732 \\
\addlinespace
V3 + CBAM + GMP          & 0.9792 & 0.9789 & 0.9822 & 0.9763 \\
\addlinespace
V4 + CBAM + GMP + Stride & 0.9767 & 0.9764 & 0.9839 & 0.9696 \\
\end{longtable}
\end{singlespace}

Three things are visible. Adding the Roboflow data raises the baseline F1 by
roughly 1.5 points, from 0.9626 to 0.9779, before any architecture change. CBAM
(V2) gives the largest single jump in precision on both sets. And V3, which adds
global max pooling, reaches the highest recall, which makes sense because max
pooling responds to sharp local defects and so catches more of the bad seeds.

\subsection{Performance on unseen data}

Validation numbers flatter a model. The real test is data the model has never
seen, drawn from a different distribution, where domain shift bites.
Tables~\ref{tab:ch6-clf-nagar-test} and \ref{tab:ch6-clf-hybrid-test} report the
same variants on an out-of-distribution test set.

\begin{singlespace}
\begin{longtable}{@{}p{0.30\textwidth}p{0.11\textwidth}p{0.10\textwidth}p{0.13\textwidth}p{0.10\textwidth}p{0.13\textwidth}@{}}
\caption{Classifier variants trained on NAGAR, evaluated on unseen data (test).}
\label{tab:ch6-clf-nagar-test}\\
\toprule
\textbf{Variant} & \textbf{Acc.} & \textbf{F1} & \textbf{Precision} & \textbf{Recall} & \textbf{Combined} \\
\midrule
\endfirsthead
\toprule
\textbf{Variant} & \textbf{Acc.} & \textbf{F1} & \textbf{Precision} & \textbf{Recall} & \textbf{Combined} \\
\midrule
\endhead
\bottomrule
\endfoot
V1 Baseline & 0.7722 & 0.5832 & 0.7568 & 0.4746 & 0.6766 \\
\addlinespace
V2 + CBAM   & 0.7920 & 0.5892 & 0.7768 & 0.4746 & 0.6906 \\
\addlinespace
V3 + GMP    & 0.7722 & 0.5162 & 0.7763 & 0.3866 & 0.6442 \\
\addlinespace
V4 + Stride & 0.7821 & 0.5349 & 0.8125 & 0.3987 & 0.6585 \\
\end{longtable}
\end{singlespace}

\begin{singlespace}
\begin{longtable}{@{}p{0.30\textwidth}p{0.11\textwidth}p{0.10\textwidth}p{0.13\textwidth}p{0.10\textwidth}p{0.13\textwidth}@{}}
\caption{Classifier variants trained on the hybrid dataset, evaluated on unseen data (test).}
\label{tab:ch6-clf-hybrid-test}\\
\toprule
\textbf{Variant} & \textbf{Acc.} & \textbf{F1} & \textbf{Precision} & \textbf{Recall} & \textbf{Combined} \\
\midrule
\endfirsthead
\toprule
\textbf{Variant} & \textbf{Acc.} & \textbf{F1} & \textbf{Precision} & \textbf{Recall} & \textbf{Combined} \\
\midrule
\endhead
\bottomrule
\endfoot
V1 Baseline & 0.8062 & 0.7477 & 0.6328 & 0.9136 & 0.7769 \\
\addlinespace
V2 + CBAM   & 0.8130 & 0.7499 & 0.6470 & 0.8918 & 0.7815 \\
\addlinespace
V3 + GMP    & 0.7725 & 0.7259 & 0.5841 & 0.9588 & 0.7492 \\
\addlinespace
V4 + Stride & 0.8318 & 0.7687 & 0.6769 & 0.8894 & 0.8003 \\
\end{longtable}
\end{singlespace}

The drop from validation to test is large and worth stating plainly. A model
trained on NAGAR alone falls from an F1 near 0.97 on validation to about 0.58 on
unseen data, because it learned the lab look of its training set. Training on the
hybrid set recovers most of that loss: the peak test F1 rises from 0.5892 to
0.7687, and recall climbs as high as 0.9588 for V3. The lesson is the one the
detection results pointed at too. Diverse training data, not a cleverer
architecture, is the main defence against domain shift. Among the hybrid-trained
variants, V4 is the most balanced, with the highest test accuracy (83.18\%) and
the best combined score (0.8003).

\subsection{The second crop}

To check that the design carries to a different seed, we ran the same four
variants on the coffee dataset. Table~\ref{tab:ch6-clf-coffee} reports the test
results.

\begin{singlespace}
\begin{longtable}{@{}p{0.30\textwidth}p{0.12\textwidth}p{0.12\textwidth}p{0.12\textwidth}p{0.12\textwidth}@{}}
\caption{Classifier variants on the coffee dataset (test).}
\label{tab:ch6-clf-coffee}\\
\toprule
\textbf{Variant} & \textbf{F1} & \textbf{Combined} & \textbf{Precision} & \textbf{Recall} \\
\midrule
\endfirsthead
\toprule
\textbf{Variant} & \textbf{F1} & \textbf{Combined} & \textbf{Precision} & \textbf{Recall} \\
\midrule
\endhead
\bottomrule
\endfoot
V1 Baseline ResNet-18    & 0.9056 & 0.8984 & 0.8898 & 0.9219 \\
\addlinespace
V2 + CBAM                & 0.9023 & 0.8957 & 0.8988 & 0.9058 \\
\addlinespace
V3 + CBAM + GMP          & 0.9101 & 0.9028 & 0.8877 & 0.9335 \\
\addlinespace
V4 + CBAM + GMP + Stride & 0.9000 & 0.8921 & 0.8803 & 0.9206 \\
\end{longtable}
\end{singlespace}

V3 is the best variant on coffee, with the highest F1 (0.9101) and recall
(0.9335), the same combination of attention and max pooling that did best on
maize. That the same design wins on two crops with different shapes is a small
piece of evidence that the attention blocks belong in the later layers of the
network, where they filter high-level features rather than low-level edges.

\section{System Performance and Latency}
\label{sec:ch6-perf}

Accuracy only matters if the system is fast enough to use. We measured the
prototype end to end on commodity hardware. A batch of three images containing
more than 200 seeds was processed in about 1.78~seconds, counting the upload, the
database writes, and the generation of the result report. That is fast enough for
a quality-control operator to work through samples without waiting on the
software.

Broken down by stage, the detector averages about 103.8~ms per image, and the
per-seed quality classifier runs in roughly 0.50 to 0.67~ms per crop, so the
classification stage is cheap even when an image contains a few hundred seeds. The
YOLOv8 fast mode brings detection down to about 5.5~ms per image for deployments
that prioritize speed. Table~\ref{tab:ch6-perf} collects these figures.

\begin{singlespace}
\begin{longtable}{@{}p{0.46\textwidth}p{0.46\textwidth}@{}}
\caption{Measured latency of the running prototype.}
\label{tab:ch6-perf}\\
\toprule
\textbf{Stage} & \textbf{Latency} \\
\midrule
\endfirsthead
\toprule
\textbf{Stage} & \textbf{Latency} \\
\midrule
\endhead
\bottomrule
\endfoot
Full batch (3 images, 200+ seeds, end to end) & $\approx$ 1.78\,s \\
\addlinespace
Detection per image (Faster R-CNN)            & $\approx$ 103.8\,ms \\
\addlinespace
Detection per image (YOLOv8 fast mode)        & $\approx$ 5.5\,ms \\
\addlinespace
Quality classification per seed               & $\approx$ 0.50--0.67\,ms \\
\end{longtable}
\end{singlespace}

\section{Limitations and Known Issues}
\label{sec:ch6-limits}

The results above also mark the limits of the work. The clearest one is the
synthetic-to-real domain gap. MultiSeedGen produces realistic composited scenes,
but a detector trained on them learns the statistics of synthetic images, and
those never match real trays exactly~\cite{tobin2017domainrand}. The generator
narrows the gap with domain-randomized lighting and degradation, but closing it
for a given crop still means validating the trained detector on genuine field
photographs before trusting it~\cite{kamilaris2018deepag}.

The classification results show the same gap from the other side. The fall from a
validation F1 near 0.97 to a test F1 near 0.58 for the NAGAR-only models is a
direct measurement of domain shift, and while the hybrid data recovers much of
it, the test scores are still well below the validation scores. Any deployment to
a new capture setting should expect to collect and label some local data and
retrain, rather than assume the validation numbers will hold.

Finally, this is a prototype rather than a finished product. It grades two crops
from top-down photographs taken in reasonable light, and it has not been hardened
or trialled in the field. The next chapter sets out what would have to change to
move it past that point.
